% ED26-BM-RAT2021-11-3: Hilgendorf, MIT2021, complete chapter pp.659--667.
\chapter{Legal Logic}
\label{sec:p11-rationality-11-3}

Legal logic in Hilgendorf's account is a theory of legal analysis and
decision making, not simply a calculus for deducing conclusions. Once an
applicable rule and its satisfied conditions are fixed, a deductive step can
be straightforward. Determining the rule's meaning, establishing the facts,
resolving conflicts of interests and deciding whether one case is relevantly
like another are not thereby settled. The chapter places legal methods along
a continuum between rule-governed decisions and weighing interests, and
argues that making their premises explicit is essential to criticism and
control \parencite[pp.~659--667]{Hilgendorf2021LegalLogic}.
\index{legal logic}\index{legal interpretation}\index{legal syllogism}

This treatment reconstructs that 2021 argument and its examples. Its principal
perspective is that of a German lawyer trained in the civil-law tradition,
with comparisons to common law. The statutory and judicial examples below
identify the source's reasoning problems; they do not state current law for
a jurisdiction or advise how a present case should be decided.

\section{Practical rationality, authority and scope}
\label{sec:p11-rat113-scope}

\begin{hbdefinition}{def:p11-rat113-legal-logic}{Legal methods as a broad evaluand}
Hilgendorf's proposed legal logic concerns the forms of argument and rules
by which lawyers conduct, or ought to conduct, legal analysis and decision
making. Its scope includes determining relevant facts, identifying and
interpreting applicable rules, subsuming a case under their conditions,
weighing protected interests and arguing from similarity.

Its \emph{formal} question is whether a conclusion follows from supplied
premises. Its \emph{interpretative} question is how a legally authoritative
text or precedent should be understood. Its \emph{evidential} question is
which facts have been established by an appropriate procedure. Its
\emph{normative and institutional} questions include which sources bind a
decision maker and which methods that decision maker ought to use.
These questions interact but are not interchangeable.
\end{hbdefinition}

The practical role of law in the source is to regulate conflicts, provide
predictability and make comparable constellations of facts subject to
comparable treatment. Legal application exercises power, so the demand for
publicly inspectable reasons is not merely a preference for elegant
arguments. Abstracting from a case's irrelevant peculiarities can constrain
arbitrary personal reactions. But what counts as legally relevant depends
on the applicable rule and its interpretation; abstraction does not certify
that the chosen rule is just or that the facts are true
\parencite[pp.~659--660]{Hilgendorf2021LegalLogic}.

Hilgendorf opposes two extremes: an image of decisions as wholly determined
by a formal legal syllogism, and free judicial determination unconstrained
by previously recognized legal requirements. His proposed middle account
retains rules while analyzing how premises are formed. He reports a tension
between lawyers' reputation for logical thought and the limited reception
of symbolic logic in legal practice, including a twentieth-century revival
in legal theory. These are the author's historical and professional
assessments, not a measured proportion of lawyers or a survey establishing
the state of legal practice today
\parencite[pp.~659--660]{Hilgendorf2021LegalLogic}.

\subsection{Civil law and common law are perspectives, not disjoint methods}

The source contrasts a civil-law emphasis on statutes, codes and regulations
with a common-law emphasis on cases and precedent. Both nevertheless require
identifying a rule and assessing relevant facts. Hilgendorf expressly notes
that precedent matters in both traditions; the distinction must not be
converted into the assertion that civil-law judges ignore previous decisions
or that common-law judges have no written law
\parencite[pp.~661--662, note~13 on p.~665]{Hilgendorf2021LegalLogic}.

In his discussion, a binding superior-court precedent and a persuasive
comparison serve different roles: the latter can bear on a case with
different material facts without mechanically determining it. Which court,
decision or reason binds in an actual jurisdiction requires its own source
and institutional analysis. Neither the formal validity of an argument nor
bare factual similarity supplies that authority
\parencite[note~28, pp.~665--666]{Hilgendorf2021LegalLogic}.

\section{What the legal syllogism establishes}
\label{sec:p11-rat113-syllogism}

The source traces the rule-to-case model through early modern legal thought,
Hobbes and the Wolff--K\"ohler tradition to Engisch's account of applying
law, first published in 1943. These are Hilgendorf's historical connections,
not a separately established priority history. Engisch treats the relevant
ought-judgments as truth-assessable statements, distinguishing them from
commands. Hilgendorf also discusses later refinements by Klug and Koch and
R\"u\ss mann; he does not identify all legal norms with ordinary descriptive
propositions \parencite[pp.~661--662]{Hilgendorf2021LegalLogic}.

\begin{hbdefinition}{def:p11-rat113-fixed-rule}{A conditional reconstruction with fixed interpretation}
Fix one interpreted legal rule, a nonempty domain of cases, an individual
case \(a\), a predicate \(T(x)\) for its complete qualifying conditions
and a predicate \(R(x)\) for the legal-consequence judgment assigned by
that rule. The narrow classical reconstruction uses
\[
                  \forall x\bigl(T(x)\to R(x)\bigr),\qquad T(a).
\]
Here \(R(a)\) is a proposition about the stipulated legal consequence,
not the physical occurrence of enforcement and not a bare imperative
given an unexplained truth value. The interpretation, authority, conditions
and factual premise are supplied inputs. This is a conditional abstraction
of the source's legal syllogism, not a complete axiomatic legal system.
\end{hbdefinition}

\begin{hblemma}{thm:p11-rat113-detachment}{proved=annex:p11.rat113-detachment}{Detachment and its premise boundary}
The premises of \cref{def:p11-rat113-fixed-rule} entail \(R(a)\)
classically. The factual premise \(T(a)\) alone does not entail \(R(a)\).
Nor do the rule and \(\neg T(a)\) entail \(\neg R(a)\).
Adding premises to an unchanged classical premise set cannot invalidate
an already valid consequence; withdrawal of a legal conclusion therefore
needs an identified change in a premise, interpretation or consequence
relation, not a failure of conditional detachment.
\end{hblemma}

\begin{hbexample}{def:p11-rat113-killing}{The source's deliberately simplified punishment example}
Hilgendorf's illustration says that qualifying killing brings a punishment
consequence; a named person is then said to have killed, and the punishment
judgment is drawn. In the reconstruction, the general rule is the major
premise, the qualifying case description the minor premise and the
case-specific consequence the conclusion. This demonstrates the inference
form. It does not establish that every killing has the same consequence:
offense classification, defenses, jurisdiction, evidence and the complete
conditions of a real provision are not contained in the simplified example.
\end{hbexample}

The internal deduction is thus much easier than the external justification
of its premises. A court can reason validly from an inapplicable rule or a
false factual description. Conversely, accepting that some applications
are defeasible does not make modus ponens invalid. If a putative exception
shows that the antecedent omitted a necessary condition, the represented
rule must be repaired; if a higher-priority norm defeats a conclusion, the
priority and defeat mechanism must be supplied separately. Hilgendorf's
chapter offers legal-methodological analysis rather than one formal default
calculus. The typed strict/defeasible distinction in
\cref{def:p18-aspic-system} is one different formal framework, not a
silent rewriting of his legal syllogism.

\section{Interpreting rules and establishing the qualifying facts}
\label{sec:p11-rat113-interpretation}

Hilgendorf decomposes a rule into a factual description and a legal
consequence linked conditionally. Actor, affected object and action can be
separate elements of the description. What matters is whether all relevant
conditions are met, not whether a case evokes an intuitively similar
conflict. His criminal-law example is theft as described through
\S~242 of the German Criminal Code in the source: movable property,
belonging to another, taking away, intent to appropriate unlawfully for
oneself or another, and the further intent requirement identified by the
author. These are distinguishable objective and subjective elements in his
analysis \parencite[pp.~660, 662]{Hilgendorf2021LegalLogic}.

\begin{hbexample}{def:p11-rat113-theft}{Decomposition does not prove subsumption}
A description that someone moved an object supplies neither ownership by
another nor the required intent in the source's theft example. Evidence
about physical movement and evidence about purpose answer different
questions. The analyst must first interpret each relevant element and then
justify its presence in the case. Merely repeating the statutory vocabulary
in a conclusion does not establish the factual or interpretative premise.
This is an explanatory contrast built from the source's elements, not an
assessment of criminal liability in a current case.
\end{hbexample}

\begin{hbmethod}{def:p11-rat113-canons}{The four interpretative methods in the source}
\begin{enumerate}
\item \emph{Wording}: investigate ordinary meaning and, where relevant,
      legal usage of the text's terms.
\item \emph{Historical interpretation}: examine legislative history and
      other relevant source material for the provision's development.
\item \emph{Systematic interpretation}: examine the provision's relation
      to other norms and their interpretations in statutes, cases and doctrine.
\item \emph{Teleological interpretation}: assess the provision's purpose
      and competing interpretations in light of their expected effects.
\end{enumerate}
Hilgendorf presents these as methods typically used successively in civil-law
analysis. They have heterogeneous sources: hermeneutic work on wording,
historical-empirical investigation, contextual linguistic and doctrinal
analysis, and evaluative reasoning informed by predictions. The list is not
a universally binding algorithm or an assurance that the four methods yield
one answer. When several meanings remain, choosing among them requires
additional reasons, including the assessment of consequences.
\end{hbmethod}

The source's social-technology interpretation concerns evaluating expected
effects of interpretations and making a reasoned selection. Forecasting an
effect and valuing it are different tasks. Neither the prediction nor the
value comparison follows solely from a word's dictionary meaning
\parencite[pp.~662--663, note~18 on p.~665]{Hilgendorf2021LegalLogic}.

Factual determination also has a procedural dimension. Hilgendorf discusses
evidence rules and logical errors as grounds for criticism or appeal, while
warning that the legal expression laws of sound thinking is less precise
than modern analyses of reasoning. His discussion of correspondence truth
is an assessment of lawyers' usual orientation, not a proof that one theory
of truth settles every evidential question. Consensus among decision makers
is not automatically correspondence with the facts. The chapter supplies no
study design or prevalence estimate for its assertions about professional
beliefs \parencite[p.~660, notes~7--8 on p.~665]{Hilgendorf2021LegalLogic}.

\section{Weighing interests and the limits of a decision-theoretic analogy}
\label{sec:p11-rat113-weighing}

Hilgendorf uses rights understood as legally protected interests to frame
conflicts about ownership, public and private planning interests, or access
to scarce medical resources. Weighing can be required by a rule rather than
being the opposite of applying any rule. Yet an invocation of weighing does
not disclose the procedure, factors or priority standard used; his central
objection is to uncontrolled substitution of personal evaluation for
reviewable rule-governed decision making
\parencite[pp.~663--664]{Hilgendorf2021LegalLogic}.
\index{weighing interests!legal reasoning}

\begin{hbexample}{def:p11-rat113-necessity}{The walking stick and the swerving motorist}
In the source's necessity illustration, a person takes a pedestrian's
walking stick to defend a child from a dog attack. The protected interest
in the child's physical integrity or life is compared with the interference
with immediate use of the stick. Hilgendorf locates this in his account of
\S~34 of the German Criminal Code, which calls for weighing the affected
interests and dangers. His example illustrates a priority comparison under
the conditions he envisages; it is not a complete current-law opinion about
every taking of property during an emergency.

In a second hypothetical case, a motorist can avoid an almost certain fatal
collision with an injured person by swerving, at a low but nonzero risk of
hitting a pedestrian. The author favors swerving within that description.
The comparison involves both the interests at stake and the probabilities
of harming them. The source supplies no numerical risks, utilities or
empirical sample, and the case is not a traffic-safety recommendation for
an actual driver.
\end{hbexample}

The source explicitly connects risk-sensitive weighing to classical
decision theory. That is a useful structural analogy, not an instruction
to put every protected interest onto a common monetary or cardinal scale.
Expected-utility choice in \cref{def:p17-expected-utility-choice} needs
its stated menu, consequence utilities and probabilities. A legal weighing
claim needs, additionally, authority to weigh these interests, admissible
factors and the relevant priority standard. The qualitative hypotheticals
do not supply those missing quantities or establish that the result follows
from a numerical maximization theorem
\parencite[p.~663, note~21 on p.~665]{Hilgendorf2021LegalLogic}.

\subsection{Life against life: lawfulness is not personal blameworthiness}

The Plank of Carneades illustrates two shipwrecked people with a plank that
will not continue to support both. The questions concern whether one may
push the other away, what could distinguish them and whether any permissible
allocation is available. Hilgendorf connects this problem to the Mignonette
maritime-cannibalism case of 1884, the switchman/trolley family, and the
German aviation-security judgment of 15 February 2006 (1 BvR 357/05)
about killing passengers to save people on the ground. These identifiers
come from the chapter's examples and notes, not an independent analysis of
the judgments. These are a problem family, not interchangeable legal
precedents or one settled calculation of the value of lives
\parencite[pp.~663--664, notes~22--24 on p.~665]{Hilgendorf2021LegalLogic}.

One could formulate a rule selecting a person with a specified characteristic,
but logic alone would not choose that characteristic. Flexible weighing
allows more circumstances to matter while also making the decision maker's
evaluations consequential. In discussing desperate life-against-life cases,
Hilgendorf distinguishes a \emph{justification}, under which an act is
lawful, from an \emph{excuse or lack of personal culpability}, under which
the act can remain unlawful while the actor is not held blameworthy.
Absence of criminal liability, by itself, does not answer which account applies.

The source's comparative discussion must retain its own qualification:
although it questions how fully the distinction has been accepted in
common-law theory, note~25 acknowledges sophisticated treatments of
justification and excuse there for decades. It would misrepresent this
discussion to claim that common law simply lacks the distinction. The
present entry adopts the conceptual contrast, not an unverified census of
national doctrines \parencite[p.~664, note~25 on p.~665]{Hilgendorf2021LegalLogic}.

Hilgendorf argues that weighing draws on historically developed practice
and a shared evaluative background, which pluralization can make less
stable. His criticism is that unclear criteria increase unpredictability
and threaten criticism and control. That is a methodological and political
argument, not evidence that every balancing decision is arbitrary or that
rules never require balancing
\parencite[p.~664, note~26 on p.~665]{Hilgendorf2021LegalLogic}.

\section{Analogy and other extensions beyond a fixed rule}
\label{sec:p11-rat113-analogy}

\begin{hbdefinition}{def:p11-rat113-analogy}{The source's substitution of a legally material element}
Let \(T_1,\ldots,T_n\), \(n\geq1\), denote interpreted conditions
of a source rule, and \(R\) its legal-consequence judgment. Put
\[
 T(x)=T_1(x)\land\bigwedge_{i=2}^n T_i(x),\qquad
 T'(x)=T'_1(x)\land\bigwedge_{i=2}^n T_i(x),
\]
with the empty conjunction true when \(n=1\). The original rule assigns
\(R(x)\) when \(T(x)\) holds. In the comparison case, \(T'_1\)
replaces \(T_1\) while the other conditions hold. The analogical question
is whether the difference is legally material, or whether there is an
adequate reason and authority to extend the consequence to \(T'\).
Similarity as an informal description does not itself define a logical
relation between these predicates.
\end{hbdefinition}

\begin{hblemma}{thm:p11-rat113-extension-boundaries}{proved=annex:p11.rat113-extension-boundaries}{Contrary and analogous conclusions require additional premises}
From \(\forall x(T(x)\to R(x))\) and \(T'(a)\), \(R(a)\)
does not follow classically. A sufficient additional premise is
\(\forall x(T'(x)\to T(x))\), if that inclusion is justified; another
is an independently justified extended rule
\(\forall x(T'(x)\to R(x))\).
Similarly, the original rule and \(\neg T(a)\) do not entail
\(\neg R(a)\). An additional exclusivity premise
\(\forall x(R(x)\to T(x))\) makes that latter inference valid.
Logical validity with these added premises does not establish their legal
authority or their substantive correctness.
\end{hblemma}

An analogical extension and a subsumption can therefore reach the same
written consequence while relying on different premises. If the new case
already satisfies the interpreted \(T_1\), applying the old rule is
subsumption; if it does not, similarity must support an extension rather
than conceal the missing condition. Hilgendorf stresses both analogy's
practical importance and the vagueness of the similarity standard
\parencite[p.~664]{Hilgendorf2021LegalLogic}.
\index{analogy!legal extension}\index{argumentum e contrario}

His criminal-law boundary is his report that analogy detrimental to an accused
is prohibited in most legal systems: precise offense conditions
limit judicial expansion of liability. It is a source-qualified
rule-of-law example, not a claim that every interpretative development is
an analogy or that one universal current prohibition has been established.
His historical warning about the removal of this constraint under Nazism
connects the issue to arbitrary punishment; note~30 is not here expanded
into a separately researched legislative history
\parencite[p.~664, note~30 on p.~666]{Hilgendorf2021LegalLogic}.

The source also places \emph{argumentum e contrario}, \emph{a fortiori}
and \emph{ad absurdum} on the range between rules and weighing. A contrary
argument needs an account of exclusivity, a stronger-case argument needs a
legally relevant ordering and transfer principle, and an absurdity argument
needs a specified unacceptable consequence. A formal reductio from
contradictory premises is different from rejecting an interpretation because
its effects are regarded as unreasonable. Hilgendorf's point is not that
classical reductio fails, but that naming one of these legal argument forms
does not supply all premises making its legal conclusion compulsory
\parencite[p.~664, note~27 on p.~665]{Hilgendorf2021LegalLogic}.

\section{Reviewable reasoning, defeasibility and automation}
\label{sec:p11-rat113-reviewability}

\begin{hbmethod}{def:p11-rat113-audit}{Making a legal-method argument inspectable}
\begin{enumerate}
\item Fix the question, jurisdictional and temporal source scope, and the
      distinction between lawfulness, liability, culpability and enforcement.
\item Identify the asserted authority and the complete interpreted rule.
      Keep an enacted text, a precedent, an interpretation and an analogical
      extension distinguishable.
\item Decompose the legally relevant conditions and state the evidence for
      each; separate factual uncertainty from uncertainty about meaning.
\item If using interpretative canons, disclose which consideration narrows
      the alternatives and which unresolved choice still needs evaluation.
\item If weighing, identify the protected interests, dangers, probabilities
      where justified, permissible factors and priority criterion. Do not
      supply invented numerical weights to make the procedure look complete.
\item If using analogy or a contrary/stronger/absurdity argument, expose the
      additional transfer, exclusion, ordering or evaluative premise.
\item Address contrary reasons and exceptions at their actual target: a
      premise, source authority, inference or priority. State the conclusion
      with its remaining conditions and uncertainty.
\end{enumerate}
This is an editorial procedure grounded in the chapter's call for transparent,
criticizable reasoning. It is not a mechanical decision algorithm or a
guarantee that every disagreement can be settled by formal logic.
\end{hbmethod}

Formal deontic logic provides a different object: for example, the standard
and dyadic obligation semantics of
\cref{def:p08-sdl,def:p08-dyadic-obligation}. Such a semantics specifies
what its formal obligations mean and which consequences follow. It does
not determine which interpretation, institution or norm is legally
authoritative. Hilgendorf's broad methodological definition should not be
identified with one of those calculi; he explicitly points to deontic logic
as a neighboring subject rather than adopting a full system
\parencite[p.~659, note~3 on p.~665]{Hilgendorf2021LegalLogic}.

Structured argumentation in \cref{sec:p18-structured-argumentation}
separately specifies argument trees, challengeable premises, defeasible
inferences, attacks and preference-conditioned defeat. It can distinguish
a challenge to a factual premise from one to an inference or conclusion;
see \cref{def:p18-aspic-attacks,def:p18-aspic-defeat}. Its explicit
framework is not merely Hilgendorf's open-ended weighing with a formal
notation attached. A legal application must justify the chosen premises,
priorities and standards of acceptance rather than infer them from the
framework's availability. His references to logical legal-argument models,
including Prakken's separate chapter~11.4, are directions for further study,
not a completeness proof for his chapter
\parencite[notes~14, 17 on p.~665]{Hilgendorf2021LegalLogic}.

The source raises automated driving as a reason to ask how traffic rules
could be represented for machine processing. It also argues, as a 2021
assessment, that the full application process is not purely deductive and
does not appear algorithmizable on then-current legal rationality standards.
This is neither a present claim about AI capabilities nor a mathematical
impossibility theorem. Encoding a fragment leaves interpretation, exceptions,
evidence and conflicting requirements to be addressed
\parencite[pp.~659--660, 662, note~4 on p.~665]{Hilgendorf2021LegalLogic}.

\begin{proofboundary}{def:p11-rat113-empirical-boundary}{A research agenda is not a measured universal result}
Hilgendorf identifies neglected work on counterfactuals, heuristics,
causation versus correlation, stereotypes, cognitive biases and probability.
His notes also suggest that legal flowcharts can be heuristics. He proposes
psychological study of weighing and analogy, and expects common socialization
to promote some similarity in decisions. The chapter does not provide a
participant population, experimental task, comparison design, effect size
or uncertainty estimate establishing that expectation.

His outlook proposes that the syllogistic structure and demand for correctly
ascertained facts have cross-cultural commonalities, and suggests anthropology
and psychology to investigate them. That proposal is not a theorem that all
legal cultures reason identically. Claims about actual behavior require the
task/population/design/evidence distinctions of
\cref{def:p19-evidence-record}; a normative requirement to establish facts
does not demonstrate that every legal system satisfies it.
\end{proofboundary}

The source's note~19 uses qadi jurisprudence as a label for the contrast
with systematic decision making. That terminology supplies no empirical
demonstration about the reasoning of traditional Islamic courts; the
methodological comparison concerns disclosed reasons versus uncontrolled
choice, not a proven classification of a legal culture.

These boundaries retain the source's distinction between structural
analysis and future empirical investigation
\parencite[pp.~660, 664--665, notes~7, 9--12, 15 on p.~665]{Hilgendorf2021LegalLogic}.

\section{Subject and task lookup}
\label{sec:p11-rat113-lookup}

Use this entry to identify the inference and premise-construction work in a
legal-method argument. A successful analysis draws the case-specific
consequence from an explicitly fixed rule and qualifying facts, while
identifying the separate interpretation and evidence justifications.
A failed analysis treats mere similarity as satisfaction of every element,
denies a consequence from failure of one sufficient condition without an
exclusivity premise, or turns an unquantified weighing into a numerical
expected-utility calculation.

For formal obligations use the named deontic semantics; for explicit
defeaters and priorities use the structured-argumentation framework; for
lottery evaluation use the finite expected-utility procedure. None of these
neighboring tasks supplies the legal authority or evaluative premises
missing from the source's methodological problem. The aim is a bounded,
reviewable account of legal reasoning, not a comprehensive law textbook or
an automated judgment in an actual dispute.
